% \iffalse meta-comment
%
% docxlike-shape.dtx 是锚在段落上的对象（线条形状与图文框）与嵌入型的图
%
% \fi
%
% \section{定位对象}
%
% \changes{v3.2a}{2026/09/26}{新增线条形状 \cs{docxshape} 与图文框 \cs{docxframe}：锚在段上，
%   纵向从段顶量，横向从栏、版心或纸的左边量；两个引擎一样，经 \cs{vadjust}~\texttt{pre}
%   放到锚定段首行之上}
% \changes{v3.2a}{2026/09/26}{新增 \cs{docxpicture}：画出来的图标成嵌入型图片，行高照装图段算}
%
% 将线条和图文框锚定到段落，并按指定参考边界定位。
%
% 相对段落的纵坐标包含本段段前间距，但不包含上段段后间距。
%
% \emph{放在哪。} 对象要画在锚定段之上，而写对象的命令在段里。\cs{vadjust}~\texttt{pre} 把
% 材料放进竖直列表、排在它所在那一行的行间胶之前，不影响行距，两个引擎都有。放进去的是一只
% 零高的盒，盒的左边就是栏的左边，与这一段怎么对齐无关。所以命令要写在锚定段的\emph{首行}
% 里，照 docx 转过来时写在段首。写在两段之间（竖直模式）的，锚在下一段上。
%
% 纵向相对页边距、页面的，要知道锚定段落在哪一页的什么位置，这一版没有做，写了报错；横向
% 相对页面的按奇数页的左边距折算。
%    \begin{macrocode}
%<*docxlike-shape>
%    \end{macrocode}
%
% \subsection{线条形状}
%
% \cs{docxshape}\oarg{键}：
% \begin{latex}
% \docxshape[
%   position = { horizontal = { absolute-position = 76.3 , to-the-right-of = column } ,
%                vertical = { absolute-position = 27.6 , below = paragraph } } ,
%   size = { width = 261.6 , height = 0 } ,
%   line = { width = 0.75 , color = auto } ]
% \end{latex}
% 键照“布局”对话框的“位置”页与“设置形状格式”窗格的“线条”：\option{absolute-position}
% 是“绝对位置”，\option{to-the-right-of} 取 \texttt{column}、\texttt{margin}、\texttt{page}，
% \option{below} 取 \texttt{paragraph}；\option{size} 是“大小”页的宽与高，线从位置那一点画到
% 宽、高之外；\option{line} 的 \option{width} 是线宽（默认 0.75 磅，与 VML 的默认一样），
% \option{color} 取 \texttt{auto} 或颜色。线的粗细居中压在坐标上。现在只有横线与竖线。
%
% VML 的 \texttt{from}、\texttt{to} 两点照写：位置取两点中靠左上的一点，宽、高取两点之差。
%    \begin{macrocode}
\dim_new:N \l__docxlike_shape_x_dim
\dim_new:N \l__docxlike_shape_y_dim
\dim_new:N \l__docxlike_shape_w_dim
\dim_new:N \l__docxlike_shape_h_dim
\dim_new:N \l__docxlike_shape_lw_dim
\tl_new:N \l__docxlike_shape_xrel_tl
\tl_new:N \l__docxlike_shape_color_tl
\msg_new:nnnn { docxlike } { shape-anchor }
  { '#1'~must~be~placed~between~paragraphs~or~in~the~first~line~of~a~paragraph. }
  { The~object~is~dropped. }
\keys_define:nn { docxlike / shape }
  {
    type .choices:nn = { line } { } ,
    position .code:n = { \__docxlike_keys_set:nn { docxlike / shape / position } {#1} } ,
    size .code:n = { \__docxlike_keys_set:nn { docxlike / shape / size } {#1} } ,
    line .code:n = { \__docxlike_keys_set:nn { docxlike / shape / line } {#1} } ,
    unknown .code:n = { \__docxlike_unknown_key: } ,
  }
\keys_define:nn { docxlike / shape / position }
  {
    horizontal .code:n = { \__docxlike_keys_set:nn { docxlike / shape / horizontal } {#1} } ,
    vertical .code:n = { \__docxlike_keys_set:nn { docxlike / shape / vertical } {#1} } ,
    unknown .code:n = { \__docxlike_unknown_key: } ,
  }
\keys_define:nn { docxlike / shape / horizontal }
  {
    absolute-position .code:n =
      { \dim_set:Nn \l__docxlike_shape_x_dim { \__docxlike_bp:n {#1} } } ,
    to-the-right-of .choices:nn = { column , margin , page }
      { \tl_set_eq:NN \l__docxlike_shape_xrel_tl \l_keys_choice_tl } ,
    unknown .code:n = { \__docxlike_unknown_key: } ,
  }
\keys_define:nn { docxlike / shape / vertical }
  {
    absolute-position .code:n =
      { \dim_set:Nn \l__docxlike_shape_y_dim { \__docxlike_bp:n {#1} } } ,
    below .choices:nn = { paragraph , line , margin , page }
      {
        \str_if_eq:VnF \l_keys_choice_tl { paragraph }
          {
            \msg_error:nne { docxlike } { not-implemented }
              { vertical~position~below~\l_keys_choice_tl }
          }
      } ,
    unknown .code:n = { \__docxlike_unknown_key: } ,
  }
\keys_define:nn { docxlike / shape / size }
  {
    width .code:n = { \dim_set:Nn \l__docxlike_shape_w_dim { \__docxlike_bp:n {#1} } } ,
    height .code:n = { \dim_set:Nn \l__docxlike_shape_h_dim { \__docxlike_bp:n {#1} } } ,
    unknown .code:n = { \__docxlike_unknown_key: } ,
  }
\keys_define:nn { docxlike / shape / line }
  {
    width .code:n = { \dim_set:Nn \l__docxlike_shape_lw_dim { \__docxlike_bp:n {#1} } } ,
    color .tl_set:N = \l__docxlike_shape_color_tl ,
    unknown .code:n = { \__docxlike_unknown_key: } ,
  }
\cs_generate_variant:Nn \msg_error:nnn { nne }
%    \end{macrocode}
%
% 横向折成相对栏的量。栏就是页边距之间（一栏的文档），相对页面的减去奇数页的左边距。
%    \begin{macrocode}
\cs_new:Npn \__docxlike_shape_page_left:
  { 1in + \hoffset + \oddsidemargin }
\cs_new_protected:Npn \__docxlike_shape_column:NV #1#2
  {
    \str_if_eq:VnT #2 { page }
      { \dim_sub:Nn #1 { \__docxlike_shape_page_left: } }
  }
%    \end{macrocode}
%
% \begin{macro}{\docxshape,\docxlike_shape:n}
%    \begin{macrocode}
\cs_new_protected:Npn \docxlike_shape:n #1
  {
    \group_begin:
      \dim_zero:N \l__docxlike_shape_x_dim
      \dim_zero:N \l__docxlike_shape_y_dim
      \dim_zero:N \l__docxlike_shape_w_dim
      \dim_zero:N \l__docxlike_shape_h_dim
      \dim_set:Nn \l__docxlike_shape_lw_dim { 0.75bp }
      \tl_set:Nn \l__docxlike_shape_xrel_tl { column }
      \tl_set:Nn \l__docxlike_shape_color_tl { auto }
      \__docxlike_keys_set:nn { docxlike / shape } {#1}
      \__docxlike_shape_column:NV \l__docxlike_shape_x_dim \l__docxlike_shape_xrel_tl
      \bool_lazy_and:nnTF
        { \dim_compare_p:nNn { \l__docxlike_shape_w_dim } > { 0pt } }
        { \dim_compare_p:nNn { \l__docxlike_shape_h_dim } > { 0pt } }
        { \msg_error:nnn { docxlike } { not-implemented } { a~slanted~line } }
        {
          \use:e
            {
              \exp_not:N \__docxlike_anchor_put:nn { \exp_not:N \docxshape }
                {
                  \exp_not:N \__docxlike_shape_line:nnnnnn
                    { \dim_use:N \l__docxlike_shape_x_dim } { \dim_use:N \l__docxlike_shape_y_dim }
                    { \dim_use:N \l__docxlike_shape_w_dim } { \dim_use:N \l__docxlike_shape_h_dim }
                    { \dim_use:N \l__docxlike_shape_lw_dim } { \l__docxlike_shape_color_tl }
                }
            }
        }
    \group_end:
  }
\NewDocumentCommand \docxshape { O { } } { \docxlike_shape:n {#1} }
%    \end{macrocode}
% \end{macro}
%
% 画线：\marg{横}\marg{纵}\marg{宽}\marg{高}\marg{线宽}\marg{颜色}，数都已折好。零高的竖盒里先往下
% 走到坐标，横线再往上退半个线宽。挂起到下一段时组已经关了，所以数是展开了传进来的。
%    \begin{macrocode}
\cs_new_protected:Npn \__docxlike_shape_line:nnnnnn #1#2#3#4#5#6
  {
    \dim_compare:nNnTF {#4} > { 0pt }
      {
        \skip_vertical:n {#2}
        \hbox_overlap_right:n
          {
            \skip_horizontal:n { #1 - #5 / 2 }
            \__docxlike_shape_color:n {#6}
            \rule {#5} {#4}
          }
      }
      {
        \skip_vertical:n { #2 - #5 / 2 }
        \hbox_overlap_right:n
          {
            \skip_horizontal:n {#1}
            \__docxlike_shape_color:n {#6}
            \rule {#3} {#5}
          }
      }
  }
\cs_new_protected:Npn \__docxlike_shape_color:n #1
  {
    \str_if_eq:nnF {#1} { auto }
      {
        \regex_match:nnTF { \A [0-9A-Fa-f]{6} \Z } {#1}
          { \color_select:nn { HTML } {#1} }
          { \color_select:n {#1} }
      }
  }
%    \end{macrocode}
%
% \subsection{放到锚定段上}
%
% \cs{\_\_docxlike\_anchor\_put:nn}\marg{命令名}\marg{画的代码}：\meta{画的代码} 在一只左边贴栏左、
% 顶在段顶的竖盒里执行。竖直模式下（两段之间）先挂起，下一段开头放；段首行里用
% \cs{vadjust}~\texttt{pre} 放到这一行之上，从那里往上退回段顶：这一段的段间距加上
% \cs{parskip}。别处报错丢掉。
%
% 盒在竖直列表里净占零：高取首行的估计高（本包撑杆的上升部与 \cs{strutbox} 的高取大），深取它的
% 相反数。不直接用零高的盒，是因为锚定段若是 \cs{vtop}（\cs{parbox}\texttt{[t]}）里的第一段，
% \cs{vtop} 按第一项的高定基线，零高的盒会把整个 \cs{vtop} 的基线挪到顶上；高取首行的高，基线仍
% 落在首行上。画的位置从盒顶算，与盒高无关。
%    \begin{macrocode}
\tl_new:N \g__docxlike_anchor_pending_tl
\cs_new_protected:Npn \__docxlike_anchor_put:nn #1#2
  {
    \mode_if_vertical:TF
      { \tl_gput_right:Nn \g__docxlike_anchor_pending_tl {#2} }
      {
        \bool_lazy_and:nnTF
          { \mode_if_horizontal_p: }
          { ! \mode_if_inner_p: }
          { \__docxlike_anchor_vadjust:n {#2} }
          { \msg_error:nnn { docxlike } { shape-anchor } { \token_to_str:N #1 } }
      }
  }
\box_new:N \l__docxlike_anchor_box
\dim_new:N \l__docxlike_anchor_dim
\cs_new_protected:Npn \__docxlike_anchor_vadjust:n #1
  {
    \tex_vadjust:D pre
      {
        \vbox_set_to_ht:Nnn \l__docxlike_anchor_box { 0pt }
          {
            \skip_vertical:n { - \l_docxlike_format_par_gap_dim - \parskip }
            #1
            \tex_vss:D
          }
        \docxlike_format_strut_height:N \l__docxlike_anchor_dim
        \dim_set:Nn \l__docxlike_anchor_dim
          { \dim_max:nn { \l__docxlike_anchor_dim } { \box_ht:N \strutbox } }
        \box_set_ht:Nn \l__docxlike_anchor_box { \l__docxlike_anchor_dim }
        \box_set_dp:Nn \l__docxlike_anchor_box { - \l__docxlike_anchor_dim }
        \box_use_drop:N \l__docxlike_anchor_box
        \tex_penalty:D 10000 \scan_stop:
      }
  }
\hook_gput_code:nnn { para / begin } { docxlike }
  {
    \tl_if_empty:NF \g__docxlike_anchor_pending_tl
      {
        \exp_args:NV \__docxlike_anchor_vadjust:n \g__docxlike_anchor_pending_tl
        \tl_gclear:N \g__docxlike_anchor_pending_tl
      }
  }
%    \end{macrocode}
%
% \subsection{图文框}
%
% \cs{docxframe}\oarg{键}\marg{内容} 是一段带 \texttt{w:framePr} 的段落，写在两段之间，锚在下一段
% 上。键照“图文框”对话框：
% \begin{latex}
% \docxframe[
%   style = Footer , text-wrapping = around ,
%   size = { width = auto } ,
%   horizontal = { position = center , relative-to = margin } ,
%   vertical = { position = 0.05 , relative-to = paragraph } ]
%   { \docxfield{PAGE} }
% \end{latex}
% \option{style} 是框里那一段的样式（默认 \texttt{Normal}），\option{format} 是这一段的直接格式
% （与 \cs{docxpar} 的可选参数一样写 \option{font}、\option{paragraph}）；\option{size} 的 \option{width} 取
% \texttt{auto}（照内容的宽，不超过栏宽）或长度，\option{height} 是最小高度；不裁剪较高的内容。
% \option{horizontal} 的 \option{position} 取 \texttt{left}、\texttt{center}、\texttt{right} 或长度，
% \option{relative-to} 取 \texttt{margin}、\texttt{page}、\texttt{column}；\option{vertical} 的
% \option{position} 取长度。普通叠放框纵向只相对 \texttt{paragraph}；LuaLaTeX 的
% \texttt{text-wrapping=around} 框也支持相对 \texttt{margin} 与 \texttt{page}。
% 矩形避让与绘制共用实际出页的锚点坐标；行盒与锚点变动时需重跑至稳定。
% \option{wrap-text} 取 \texttt{largest}（默认）、\texttt{left-only}、\texttt{right-only}、\texttt{both-sides}。
% 两侧模式把同一行的行片合并到共同基线，孤行与整段约束按物理行算。没有足够侧边空间时，
% 先避让到遮挡下缘再排正文；图片还支持 \texttt{text-wrapping=top-and-bottom}。
% 轮廓绕排、多栏与向锚定段之前的正文回溯绕排尚不支持。
% 水平/垂直子键的 \option{distance-from-text} 设置左右/上下相同间距；顶层同名键
% 可按 \texttt{left/right/top/bottom} 设置不同的非负间距，默认零。
%
% 框里那一段照样式排，段前段后算在框里。排之前把格式模块的全局状态存下、排完放回，同页眉
% 页脚那样（\cs{\_\_docxlike\_headfoot\_state\_save:}），另加上一段的段后与段间距那几个量：框
% 不占正文流，下一段的段前照旧对着框前面那一段算。
%    \begin{macrocode}
\tl_new:N \l__docxlike_frame_style_tl
\tl_new:N \l__docxlike_frame_format_tl
\tl_new:N \l__docxlike_frame_width_tl
\tl_new:N \l__docxlike_frame_height_tl
\tl_new:N \l__docxlike_frame_xpos_tl
\tl_new:N \l__docxlike_frame_xrel_tl
\tl_new:N \l__docxlike_frame_yrel_tl
\tl_new:N \l__docxlike_frame_wrap_tl
\tl_new:N \l__docxlike_frame_side_tl
\dim_new:N \l__docxlike_frame_y_dim
\dim_new:N \l__docxlike_frame_w_dim
\dim_new:N \l__docxlike_frame_left_dim
\dim_new:N \l__docxlike_frame_right_dim
\dim_new:N \l__docxlike_frame_top_dim
\dim_new:N \l__docxlike_frame_bottom_dim
\int_new:N \g__docxlike_frame_int
\int_new:N \g__docxlike_wrap_int
\msg_new:nnn { docxlike } { rectangular-wrap-context }
  { Rectangular~wrapping~requires~LuaLaTeX,~Word~linebreaking,~and~an~object~between~single-column~body~paragraphs. }
\msg_new:nnn { docxlike } { rectangular-wrap-distance }
  { Distances~from~text~must~be~nonnegative. }
\msg_new:nnn { docxlike } { picture-inline-options }
  { Position~and~distance~options~require~floating~text-wrapping. }
\keys_define:nn { docxlike / frame }
  {
    style .tl_set:N = \l__docxlike_frame_style_tl ,
    format .tl_set:N = \l__docxlike_frame_format_tl ,
    text-wrapping .choices:nn = { none , around }
      { \tl_set_eq:NN \l__docxlike_frame_wrap_tl \l_keys_choice_tl } ,
    wrap-text .choices:nn = { largest , left-only , right-only , both-sides }
      { \tl_set_eq:NN \l__docxlike_frame_side_tl \l_keys_choice_tl } ,
    distance-from-text .code:n = { \__docxlike_keys_set:nn { docxlike / frame / distance } {#1} } ,
    size .code:n = { \__docxlike_keys_set:nn { docxlike / frame / size } {#1} } ,
    horizontal .code:n = { \__docxlike_keys_set:nn { docxlike / frame / horizontal } {#1} } ,
    vertical .code:n = { \__docxlike_keys_set:nn { docxlike / frame / vertical } {#1} } ,
    unknown .code:n = { \__docxlike_unknown_key: } ,
  }
\keys_define:nn { docxlike / frame / size }
  {
    width .tl_set:N = \l__docxlike_frame_width_tl ,
    height .tl_set:N = \l__docxlike_frame_height_tl ,
    unknown .code:n = { \__docxlike_unknown_key: } ,
  }
\keys_define:nn { docxlike / frame / horizontal }
  {
    position .tl_set:N = \l__docxlike_frame_xpos_tl ,
    relative-to .choices:nn = { margin , page , column }
      { \tl_set_eq:NN \l__docxlike_frame_xrel_tl \l_keys_choice_tl } ,
    distance-from-text .code:n =
      {
        \dim_set:Nn \l__docxlike_frame_left_dim { \__docxlike_bp:n {#1} }
        \dim_set_eq:NN \l__docxlike_frame_right_dim \l__docxlike_frame_left_dim
      } ,
    unknown .code:n = { \__docxlike_unknown_key: } ,
  }
\keys_define:nn { docxlike / frame / vertical }
  {
    position .code:n =
      {
        \str_case:nnTF {#1} { { top } { } { center } { } { bottom } { } }
          { \msg_error:nnn { docxlike } { not-implemented } { vertical~position~#1 } }
          { \dim_set:Nn \l__docxlike_frame_y_dim { \__docxlike_bp:n {#1} } }
      } ,
    relative-to .choices:nn = { paragraph , margin , page }
      { \tl_set_eq:NN \l__docxlike_frame_yrel_tl \l_keys_choice_tl } ,
    distance-from-text .code:n =
      {
        \dim_set:Nn \l__docxlike_frame_top_dim { \__docxlike_bp:n {#1} }
        \dim_set_eq:NN \l__docxlike_frame_bottom_dim \l__docxlike_frame_top_dim
      } ,
    unknown .code:n = { \__docxlike_unknown_key: } ,
  }
\keys_define:nn { docxlike / frame / distance }
  {
    left .code:n = { \dim_set:Nn \l__docxlike_frame_left_dim { \__docxlike_bp:n {#1} } } ,
    right .code:n = { \dim_set:Nn \l__docxlike_frame_right_dim { \__docxlike_bp:n {#1} } } ,
    top .code:n = { \dim_set:Nn \l__docxlike_frame_top_dim { \__docxlike_bp:n {#1} } } ,
    bottom .code:n = { \dim_set:Nn \l__docxlike_frame_bottom_dim { \__docxlike_bp:n {#1} } } ,
    unknown .code:n = { \__docxlike_unknown_key: } ,
  }
\cs_new_protected:Npn \__docxlike_frame_defaults:
  {
    \tl_set:Nn \l__docxlike_frame_style_tl { Normal }
    \tl_clear:N \l__docxlike_frame_format_tl
    \tl_set:Nn \l__docxlike_frame_width_tl { auto }
    \tl_set:Nn \l__docxlike_frame_height_tl { auto }
    \tl_set:Nn \l__docxlike_frame_xpos_tl { left }
    \tl_set:Nn \l__docxlike_frame_xrel_tl { column }
    \tl_set:Nn \l__docxlike_frame_yrel_tl { paragraph }
    \tl_set:Nn \l__docxlike_frame_wrap_tl { none }
    \tl_set:Nn \l__docxlike_frame_side_tl { largest }
    \dim_zero:N \l__docxlike_frame_y_dim
    \dim_zero:N \l__docxlike_frame_left_dim
    \dim_zero:N \l__docxlike_frame_right_dim
    \dim_zero:N \l__docxlike_frame_top_dim
    \dim_zero:N \l__docxlike_frame_bottom_dim
  }
\cs_new_protected:Npn \__docxlike_wrap_put:
  {
    \dim_compare:nNnT
      { \dim_min:nn { \l__docxlike_frame_left_dim }
          { \dim_min:nn { \l__docxlike_frame_right_dim }
              { \dim_min:nn { \l__docxlike_frame_top_dim } { \l__docxlike_frame_bottom_dim } } } }
      < { 0pt } { \msg_error:nn { docxlike } { rectangular-wrap-distance } }
    \bool_lazy_all:nTF
      {
        { \sys_if_engine_luatex_p: }
        { \mode_if_vertical_p: }
        { ! \mode_if_inner_p: }
        { \bool_if_p:N \g__docxlike_linebreak_word_bool }
        { \bool_if_p:N \g_docxlike_word_layout_bool }
        { ! \bool_if_p:N \g__docxlike_cols_open_bool }
      }
      {
        \int_gincr:N \g__docxlike_wrap_int
        \lua_now:e
          {
            docxlike.linebreak.wrap_add(\int_use:N \g__docxlike_wrap_int,~{
              x=\dim_to_decimal_in_sp:n { \l__docxlike_shape_x_dim },
              y=\dim_to_decimal_in_sp:n { \l__docxlike_frame_y_dim },
              w=\dim_to_decimal_in_sp:n { \l__docxlike_frame_w_dim },
              h=\dim_to_decimal_in_sp:n
                { \box_ht:c { g__docxlike_frame_\int_use:N \g__docxlike_frame_int _box }
                  + \box_dp:c { g__docxlike_frame_\int_use:N \g__docxlike_frame_int _box } },
              left=\dim_to_decimal_in_sp:n { \l__docxlike_frame_left_dim },
              right=\dim_to_decimal_in_sp:n { \l__docxlike_frame_right_dim },
              top=\dim_to_decimal_in_sp:n { \l__docxlike_frame_top_dim },
              bottom=\dim_to_decimal_in_sp:n { \l__docxlike_frame_bottom_dim },
              side="\l__docxlike_frame_side_tl",
              wrap="\l__docxlike_frame_wrap_tl",
              xalign="\l__docxlike_frame_xpos_tl",
              xrel="\l__docxlike_frame_xrel_tl",~yrel="\l__docxlike_frame_yrel_tl",
              page_left=\dim_to_decimal_in_sp:n { \__docxlike_shape_page_left: },
              page_top=\dim_to_decimal_in_sp:n { 1in+\voffset+\topmargin+\headheight+\headsep }
            })
          }
        \use:e
          {
            \exp_not:N \__docxlike_anchor_put:nn { \exp_not:N \docxframe }
              {
                \exp_not:N \__docxlike_wrap_draw:nn
                  { \int_use:N \g__docxlike_wrap_int } { \int_use:N \g__docxlike_frame_int }
              }
          }
      }
      { \msg_error:nn { docxlike } { rectangular-wrap-context } }
  }
\box_new:N \l__docxlike_wrap_draw_box
\cs_new_protected:Npn \__docxlike_wrap_draw:nn #1#2
  {
    \vbox_set_to_ht:Nnn \l__docxlike_wrap_draw_box { 0pt }
      {
        \lua_shipout:n { docxlike.linebreak.wrap_record("f"~..~#1,~0,~0) }
        \lua_now:n { docxlike.linebreak.wrap_place(#1) }
        \skip_vertical:n { \l__docxlike_frame_y_dim }
        \hbox_overlap_right:n
          {
            \skip_horizontal:n { \l__docxlike_shape_x_dim }
            \box_use_drop:c { g__docxlike_frame_#2_box }
          }
        \tex_vss:D
      }
    \box_set_dp:Nn \l__docxlike_wrap_draw_box { 0pt }
    \tex_prevdepth:D = -1000pt \scan_stop:
    \box_use_drop:N \l__docxlike_wrap_draw_box
  }
\sys_if_engine_luatex:T
  {
    \hook_gput_code:nnn { shipout / before } { docxlike-wrap }
      { \lua_now:n { if~docxlike~and~docxlike.linebreak~then~docxlike.linebreak.wrap_page=docxlike.linebreak.wrap_page+1~end } }
    \hook_gput_code:nnn { enddocument / afterlastpage } { docxlike-wrap }
      { \lua_now:n { if~docxlike~and~docxlike.linebreak~then~docxlike.linebreak.wrap_finish()~end } }
  }
%    \end{macrocode}
%
% \begin{macro}{\docxframe,\docxlike_frame:nn}
% 宽取 \texttt{auto} 时先把内容照样式的字体排成一行量宽。框里的盒按序号存，挂起的代码里只记
% 序号与坐标。
%    \begin{macrocode}
\cs_new_protected:Npn \docxlike_frame:nn #1#2
  {
    \group_begin:
      \__docxlike_frame_defaults:
      \__docxlike_keys_set:nn { docxlike / frame } {#1}
      \int_gincr:N \g__docxlike_frame_int
      \box_if_exist:cF { g__docxlike_frame_ \int_use:N \g__docxlike_frame_int _box }
        { \box_new:c { g__docxlike_frame_ \int_use:N \g__docxlike_frame_int _box } }
      \__docxlike_headfoot_state_save:
      \__docxlike_frame_state_save:
      \str_if_eq:VnTF \l__docxlike_frame_width_tl { auto }
        {
          \hbox_set:Nn \l__docxlike_headfoot_box
            {
              \docxlike_format_with:VVn \l__docxlike_frame_style_tl \l__docxlike_frame_format_tl
                { \docxlike_format_apply_state_else:Vn \l__docxlike_frame_style_tl { } }
              #2
            }
          \dim_set:Nn \l__docxlike_frame_w_dim
            { \dim_min:nn { \box_wd:N \l__docxlike_headfoot_box } { \linewidth } }
        }
        { \dim_set:Nn \l__docxlike_frame_w_dim { \__docxlike_bp:V \l__docxlike_frame_width_tl } }
      \vbox_gset:cn { g__docxlike_frame_ \int_use:N \g__docxlike_frame_int _box }
        {
          \dim_set_eq:NN \hsize \l__docxlike_frame_w_dim
          \dim_set_eq:NN \linewidth \l__docxlike_frame_w_dim
          \dim_zero:N \parskip
          \dim_zero:N \parindent
          \everypar { }
          \docxlike_format_par:VVn \l__docxlike_frame_style_tl \l__docxlike_frame_format_tl {#2}
        }
      \__docxlike_frame_state_restore:
      \__docxlike_headfoot_state_restore:
      \str_if_eq:VnF \l__docxlike_frame_height_tl { auto }
        {
          \box_gset_ht:cn { g__docxlike_frame_\int_use:N \g__docxlike_frame_int _box }
            {
              \dim_max:nn
                { \box_ht:c { g__docxlike_frame_\int_use:N \g__docxlike_frame_int _box } }
                { \__docxlike_bp:V \l__docxlike_frame_height_tl
                  - \box_dp:c { g__docxlike_frame_\int_use:N \g__docxlike_frame_int _box } }
            }
        }
      \__docxlike_frame_x:
      \str_if_eq:VnTF \l__docxlike_frame_wrap_tl { around }
        { \__docxlike_wrap_put: }
        {
          \str_if_eq:VnF \l__docxlike_frame_yrel_tl { paragraph }
            { \msg_error:nne { docxlike } { not-implemented } { unwrapped~frame~vertical~position~relative~to~\l__docxlike_frame_yrel_tl } }
          \use:e
            {
              \exp_not:N \__docxlike_anchor_put:nn { \exp_not:N \docxframe }
                {
                  \exp_not:N \skip_vertical:n { \dim_use:N \l__docxlike_frame_y_dim }
                  \exp_not:N \hbox_overlap_right:n
                    {
                      \exp_not:N \skip_horizontal:n { \dim_use:N \l__docxlike_shape_x_dim }
                      \exp_not:N \box_use_drop:c
                        { g__docxlike_frame_ \int_use:N \g__docxlike_frame_int _box }
                    }
                }
            }
        }
    \group_end:
  }
\cs_generate_variant:Nn \docxlike_format_apply_state_else:nn { V }
\dim_new:N \l__docxlike_frame_prev_after_dim
\dim_new:N \l__docxlike_frame_pending_after_dim
\dim_new:N \l__docxlike_frame_gap_dim
\bool_new:N \l__docxlike_frame_placed_bool
\bool_new:N \l__docxlike_frame_gap_bool
\tl_new:N \l__docxlike_frame_anchor_pending_tl
\cs_new_protected:Npn \__docxlike_frame_state_save:
  {
    \tl_set_eq:NN \l__docxlike_frame_anchor_pending_tl \g__docxlike_anchor_pending_tl
    \tl_gclear:N \g__docxlike_anchor_pending_tl
    \dim_set_eq:NN \l__docxlike_frame_prev_after_dim \g__docxlike_format_prev_after_dim
    \dim_set_eq:NN \l__docxlike_frame_pending_after_dim \g__docxlike_format_pending_after_dim
    \dim_set_eq:NN \l__docxlike_frame_gap_dim \g__docxlike_format_gap_dim
    \bool_set_eq:NN \l__docxlike_frame_placed_bool \g__docxlike_format_after_placed_bool
    \bool_set_eq:NN \l__docxlike_frame_gap_bool \g__docxlike_format_gap_bool
  }
\cs_new_protected:Npn \__docxlike_frame_state_restore:
  {
    \tl_gset_eq:NN \g__docxlike_anchor_pending_tl \l__docxlike_frame_anchor_pending_tl
    \dim_gset_eq:NN \g__docxlike_format_prev_after_dim \l__docxlike_frame_prev_after_dim
    \dim_gset_eq:NN \g__docxlike_format_pending_after_dim \l__docxlike_frame_pending_after_dim
    \dim_gset_eq:NN \g__docxlike_format_gap_dim \l__docxlike_frame_gap_dim
    \bool_gset_eq:NN \g__docxlike_format_after_placed_bool \l__docxlike_frame_placed_bool
    \bool_gset_eq:NN \g__docxlike_format_gap_bool \l__docxlike_frame_gap_bool
  }
\cs_generate_variant:Nn \docxlike_format_par:nnn { VV }
\cs_generate_variant:Nn \docxlike_format_with:nnn { VV }
\cs_generate_variant:Nn \__docxlike_bp:n { V }
\NewDocumentCommand \docxframe { O { } +m } { \docxlike_frame:nn {#1} {#2} }
%    \end{macrocode}
% \end{macro}
%
% 横向：左、中、右按栏宽（相对页面时按纸宽）与框宽算，长度照写，再折成相对栏的量。
%    \begin{macrocode}
\cs_new_protected:Npn \__docxlike_frame_x:
  {
    \dim_set:Nn \l__docxlike_shape_x_dim
      {
        \str_case:VnF \l__docxlike_frame_xpos_tl
          {
            { left } { 0pt }
            { center } { ( \__docxlike_frame_span: - \l__docxlike_frame_w_dim ) / 2 }
            { right } { \__docxlike_frame_span: - \l__docxlike_frame_w_dim }
          }
          { \__docxlike_bp:V \l__docxlike_frame_xpos_tl }
      }
    \__docxlike_shape_column:NV \l__docxlike_shape_x_dim \l__docxlike_frame_xrel_tl
  }
\cs_new:Npn \__docxlike_frame_span:
  { \str_if_eq:VnTF \l__docxlike_frame_xrel_tl { page } { \paperwidth } { \linewidth } }
\cs_generate_variant:Nn \str_case:nnF { VnF }
%    \end{macrocode}
%
% \subsection{嵌入型的图}
%
% \begin{macro}{\docxpicture,\docxlike_picture:n}
% 环绕方式是“嵌入型”的图随正文走，行高照装图段的规矩定：不贴行网格时上升部是图高，多倍行距
% 多出来的那一截落在图下；贴格时占整数个网格行，图在里面居中。\cs{includegraphics} 出的图
% 断行层认得出来。用 PDF literal 或 TikZ 画在盒子里的图认不出，当成一个很高的字，行高就照字排了。
% \cs{docxpicture}\marg{内容} 把内容装进一只横盒、标上属性，断行层见到这个属性就照图排，盒里
% 有字也不管。只有 \hologo{LuaLaTeX} 上有断行层，\hologo{XeLaTeX} 上只装盒。
%    \begin{macrocode}
\sys_if_engine_luatex:T { \newattribute \docxlike@picture@attr }
\cs_new_protected:Npn \docxlike_picture:n #1
  {
    \mode_leave_vertical:
    \group_begin:
      \sys_if_engine_luatex:T { \docxlike@picture@attr = 1 \scan_stop: }
      \hbox:n {#1}
    \group_end:
  }
\bool_new:N \l__docxlike_picture_position_bool
\keys_define:nn { docxlike / picture }
  {
    text-wrapping .choices:nn = { inline , square , top-and-bottom }
      { \tl_set_eq:NN \l__docxlike_frame_wrap_tl \l_keys_choice_tl } ,
    wrap-text .code:n =
      { \bool_set_true:N \l__docxlike_picture_position_bool \__docxlike_keys_set:nn { docxlike / frame } { wrap-text={#1} } } ,
    horizontal .code:n =
      { \bool_set_true:N \l__docxlike_picture_position_bool \__docxlike_keys_set:nn { docxlike / frame / horizontal } {#1} } ,
    vertical .code:n =
      { \bool_set_true:N \l__docxlike_picture_position_bool \__docxlike_keys_set:nn { docxlike / frame / vertical } {#1} } ,
    distance-from-text .code:n =
      { \bool_set_true:N \l__docxlike_picture_position_bool \__docxlike_keys_set:nn { docxlike / frame / distance } {#1} } ,
    unknown .code:n = { \__docxlike_unknown_key: } ,
  }
\cs_new_protected:Npn \docxlike_picture:nn #1#2
  {
    \group_begin:
      \__docxlike_frame_defaults:
      \bool_set_false:N \l__docxlike_picture_position_bool
      \tl_set:Nn \l__docxlike_frame_wrap_tl { inline }
      \__docxlike_keys_set:nn { docxlike / picture } {#1}
      \str_if_eq:VnTF \l__docxlike_frame_wrap_tl { inline }
        {
          \bool_if:NT \l__docxlike_picture_position_bool
            { \msg_error:nn { docxlike } { picture-inline-options } }
          \docxlike_picture:n {#2}
        }
        {
          \int_gincr:N \g__docxlike_frame_int
          \box_new:c { g__docxlike_frame_\int_use:N \g__docxlike_frame_int _box }
          \__docxlike_headfoot_state_save:
          \__docxlike_frame_state_save:
          \hbox_gset:cn { g__docxlike_frame_\int_use:N \g__docxlike_frame_int _box } {#2}
          \__docxlike_frame_state_restore:
          \__docxlike_headfoot_state_restore:
          \dim_set:Nn \l__docxlike_frame_w_dim
            { \box_wd:c { g__docxlike_frame_\int_use:N \g__docxlike_frame_int _box } }
          \__docxlike_frame_x:
          \__docxlike_wrap_put:
        }
    \group_end:
  }
\NewDocumentCommand \docxpicture { O { } +m } { \docxlike_picture:nn {#1} {#2} }
%    \end{macrocode}
% \end{macro}
%
%    \begin{macrocode}
%</docxlike-shape>
%    \end{macrocode}
%
% \Finale
